\section{Related Works}

\paragraph{Muon and matrix-aware optimizers.}
A growing line of work argues that modern neural networks should not be optimized solely with coordinate-wise Euclidean methods, because their dominant parameters are matrices whose geometry is better captured at the operator level. Classical examples include curvature- or matrix-aware methods such as K-FAC, Shampoo, Adafactor, and spectral or matrix-preconditioned optimizers, which exploit layerwise structure rather than treating every entry independently~\cite{martens2015kfac,gupta2018shampoo,shazeer2018adafactor,carlson2015psd,carlson2015ssd,shi2023distributedshampoo}. Recent interest has further shifted from matrix preconditioning to matrix \emph{orthogonalization}: Muon replaces the raw gradient by a polarized direction based on the matrix sign operator, implemented either through exact SVD or fast Newton--Schulz-type iterations~\cite{polar_express_2025,kovalev2025gradientorth,bernstein2024oldoptimizer,he2025lowrank,liu2025muonscalable,shah2025practicalefficiency}. This family has stimulated both theoretical and algorithmic follow-up work, including analyses under spectral-norm constraints, convergence studies, low-rank or polynomial approximations, and unifying views through linear minimization or preconditioning perspectives~\cite{chen2025spectralnorm,li2025note,lau2025polargrad,sfyraki2025lionsmuons,pethick2025normconstrained,grishina2025chebyshev}. These results substantially clarify why orthogonalized updates can be effective for training large models, yet they are almost exclusively developed for centralized or data-parallel settings where the matrix-polarization step is not entangled with decentralized consensus.

\paragraph{Decentralized optimization and heterogeneity correction.}
On the systems side, decentralized optimization has long been studied as a serverless alternative to parameter-server training. Early decentralized SGD and gossip-style methods demonstrated that peer-to-peer learning can remove central bottlenecks and even achieve attractive communication/computation trade-offs, but they also revealed a central challenge: under data heterogeneity, local iterates drift apart and consensus errors bias optimization~\cite{lian2017dpsgd,lian2017async,tang2018d2}. A large literature addresses this issue via exact correction, primal--dual reformulations, or gradient tracking. Canonical milestones include EXTRA and its proximal variants~\cite{shi2015extra,shi2015pgextra}, DIGing and related tracking-based methods~\cite{nedic2017diging,pu2021dsgt,zhang2019dsgtnonconvex}, and exact-diffusion / D$^2$-style methods that compensate heterogeneity-induced bias while preserving decentralized communication~\cite{yuan2019exactdiffusion,tang2018d2}. More recent analyses refined this picture for stochastic and nonconvex learning, establishing topology-separated rates, transient-time improvements, and linear-speedup characterizations~\cite{alghunaim2022unified,xin2021topologyindependent,yuan2021removing,lu2020optimal,koloskova2020unified,hhuang2021transient,haomu2023speedup}. In parallel, unified single-loop primal--dual viewpoints have started to systematize seemingly different correction mechanisms and expose common polynomial-in-mixing-matrix structures behind EXTRA-, exact-diffusion-, and tracking-type updates~\cite{zhu2024sparkle}. However, these frameworks are fundamentally derived for Euclidean gradients or proximal maps; they do not directly model the nonlinear matrix orthogonalization step that defines Muon.

\paragraph{Toward decentralized Muon.}
Only very recent works begin to combine matrix-aware optimization with decentralized communication. DeMuon is, to our knowledge, the first explicit attempt to extend Muon to graph-based decentralized learning by coupling Muon's Newton--Schulz orthogonalization with gradient tracking~\cite{he2025demuon}. Closely related distributed orthonormalized-update methods such as Dion, as well as decentralized matrix problems like DeEPCA, also suggest that non-Euclidean matrix operations can be implemented over networks~\cite{ahn2025dion,ye2021deepca}. These papers are important first steps, but they also expose the current limitations of the area. First, existing decentralized Muon designs are tied to a \emph{single} communication template---typically gradient tracking---rather than a unified recursion that can also recover exact-diffusion or EXTRA-like corrections. Second, current analyses do not leverage the refined topology-separation phenomena developed in unified decentralized theory, so it remains unclear whether matrix orthogonalization can enjoy the same higher-order network dependence observed in SUDA-style analyses~\cite{alghunaim2022unified,zhu2024sparkle,haomu2023speedup}. Third, although asymptotic linear speedup is a defining benchmark for decentralized SGD~\cite{lian2017dpsgd,lu2020optimal,koloskova2020unified}, the interaction between local nonlinear polarization and network averaging is still poorly understood. In particular, whether serverless Muon can simultaneously exploit matrix geometry, correct heterogeneity, and retain decentralized linear speedup remains open. Our work is motivated precisely by these three gaps: a unified algorithmic view beyond gradient tracking, a sharper theory that separates topology effects, and a speedup analysis tailored to matrix-aware decentralized optimization.
